%%%% CAPÍTULO 2 - REVISÃO DA LITERATURA (OU REVISÃO BIBLIOGRÁFICA, ESTADO DA ARTE, ESTADO DO CONHECIMENTO)
%%
%% O autor deve registrar seu conhecimento sobre a literatura básica do assunto, discutindo e comentando a informação já publicada.
%% A revisão deve ser apresentada, preferencialmente, em ordem cronológica e por blocos de assunto, procurando mostrar a evolução do tema.
%% Título e rótulo de capítulo (rótulos não devem conter caracteres especiais, acentuados ou cedilha)
\chapter{Referencial te\'orico}\label{cap:referencialTeorico}

%% ESCREVER UMA INTRODUÇÃO PARA ESSE CAPÍTULO

\section{Teoria de Grafos}\label{sec:teoria_grafos}

%% ESCREVER UMA INTRODUÇÃO PARA ESSE SEÇÃO

\subsection{Coeficiente de \textit{Clustering} Local}\label{subsec:clustering_local}

A existência de agrupamentos ou \textit{clusters} em redes complexas foi o que motivou a criação de
uma variedade de métricas com a intenção de quantificar sua prevalência, sendo uma delas o
coeficiente de clustering \cite{Barabasi69-70}. Existem duas variações dessa medida, a global e a local.

O coeficiente de \textit{clustering} global quantifica o agrupamento geral de um grafo em termos do
número de triângulos existentes. Ele se baseia em conjuntos de três vértices chamados
trio. Um trio é composto por três vértices conectados por duas – chamado trio aberto – ou três
- chamado trio fechado – arestas não direcionadas. Um triângulo consiste em três trios
fechados, um em cada vértice. Portanto, o coeficiente de clustering global calcula a razão
entre o número total de trios fechados e o número total de trios, abertos ou fechados \cite[48--49]{Barabasi69-70}. Porém, se o objetivo é analisar algumas características mais específicas como modularidade,
estrutura de comunidade, entre outras, o conceito de coeficiente de \textit{clustering} local se adequa
melhor a necessidade \cite[63]{Lima2022}. Assim, como o objetivo desse trabalho está focado nessa outra
medida, ela será abordada de forma mais profunda.

O coeficiente de \textit{clustering} local é uma medida importante que quantifica a tendência de
formação de grupos (ou \textit{clusters}) entre os nós de um grafo. Ela é uma métrica que avalia a
densidade de conexões entre os vizinhos de um nó específico $v_i$ em um grafo, e fornece uma
medida da probabilidade de que dois vizinhos desse nó estejam também conectados entre si,
formando um triângulo no grafo. Ao analisar um grafo que modela uma rede social, por
exemplo, onde cada nó simboliza uma pessoa e cada aresta representa um relacionamento
entre duas pessoas, o coeficiente de \textit{clustering} local pode ser utilizado para medir a
probabilidade de que uma pessoa forme comunidades dentro dessa rede. Isso é feito estimando
a chance de que dois amigos de uma pessoa também se tornem amigos entre si \cite[11]{Lima2022}. Esta
medida foi originalmente proposta por Watts e Strogatz em 1998 para determinar se um grafo
constitui uma rede de mundo pequeno.

Um grafo $G = (V, E)$ é composto, formalmente, por um conjunto de vértices $V$ e de um
conjunto de arestas $E$, onde a aresta $e_{ij}$ conecta o vértice $v_i$ ao vértice $v_j$. A vizinhança de $N_i$
para um vértice $v_i$ inclui todos os seus vértices adjacentes e as arestas que os conectam, ou
seja, a vizinhança abrange todos os nós diretamente ligados. E definimos $k_i$ como sendo o
número de vértices de $N_i$. \cite[23--24]{Easley2010}. Em termos matemáticos, temos que:

\begin{equation}
    N_i = \{ v_i : e_{ij} \in E \land e_{ji} \in E\}
\end{equation}

\begin{equation}
    k_i = |N_i|
\end{equation}

A definição formal do coeficiente de \textit{clustering} local é definida pela razão entre o número de arestas existentes entre os vizinhos de $v_i$ e o número total de arestas possíveis entre esses
vizinhos. O número de arestas existentes entre os vizinhos de $v_i$ é dado por $E_i$. O número
máximo de arestas possíveis entre os $k_i$ vizinhos de $v_i$ é:

\begin{equation}
    \frac{k_i(k_i-1)}{2}
\end{equation}

Sendo assim, o coeficiente de \textit{clustering} local $C_i$ é então definido como:

\begin{equation}
    C_i = \frac{2E_i}{k_i(k_i-1)}
\end{equation}

Este coeficiente varia entre 0 e 1, onde 0 indica que nenhum dos vizinhos de $v_i$ estão
conectados entre si e 1 indica que todos os vizinhos de $v_i$ estão completamente conectados,
formando um clique \cite[49]{Barabasi69-70}.

Abaixo está uma figura exemplificando o cálculo para o vértice em azul, onde as arestas em
destaque preto são as arestas existentes entre seus vizinhos e as arestas tracejadas em
vermelhos são as possíveis arestas entre os vizinhos.

%%COLOCAR IMAGEM AQUI DEPOIS - FAZER UMA DE AUTORIA PROPRIA

Um valor alto de $C_i$ sugere uma forte coesão local, onde os vizinhos de $v_i$ tendem a formar
um subgrafo densamente conectado. Isso é comum em redes sociais, onde amigos de um
indivíduo tendem a ser também amigos entre si, por exemplo.

\section{Algoritmo de Aproximação}\label{sec:algoritmo_aprox}

Existe uma variedade de algoritmos já propostos que buscam computar o coeficiente de
\textit{clustering} local de um grafo, tanto tendo abordagens para cálculos exatos quanto para cálculos aproximados.

Os algoritmos exatos propostos para calcular o coeficiente de clustering local de cada vértice
de um grafo estão estreitamente relacionados ao problema de contagem de triângulos locais,
já que este último pode ser transformado no primeiro \cite[64]{Lima2022}. Um algoritmo de “força bruta” que lista todos os trios de vértices para verificar quantos deles possui triângulos executa em tempo cúbico. O algoritmo proposto por Alon et al. (1997) %% citar
melhora esse tempo de execução para $O(m^{\frac{2w}{(w+1)}})$, onde $m = |E|$ e $w$ é 2,373, o expoente da multiplicação de matrizes \cite[209--223]{Alon1997}. Já o algoritmo proposto por Chiba e Nishizeki et al. (1985) %%citar 
melhora esse tempo de execução para $O(\alpha(G) \cdot m)$, onde $\alpha(G)$ é a arborescência do grafo $G$. Como $\alpha(G) = O(\sqrt{m})$, esse algoritmo executa em $O(m^(\frac{3}{2}))$ no pior caso \cite[210--223]{Chiba1985}.
Porém, tendo o cenário que a entrada é um grafo que modela uma rede complexa, onde o grau
é na casa dos milhares, utilizar esses algoritmos exatos se torna computacionalmente inviável.
Logo, algumas abordagens de aproximação que são mais rápidas se fazem mais interessantes.
Uma dessas abordagens é o algoritmo proposto por Kutzkov e Pagh em 2013, onde os autores apresentam um algoritmo de streaming com $\epsilon$-aproximação para estimar, com probabilidade de $1-\delta$, o coeficiente de \textit{clustering} local de cada vértice com grau pelo menos $d$ em tempo esperado $O(\frac{m}{\alpha\epsilon^2}log\frac{1}{\epsilon}log\frac{n}{\delta})$, onde $n = |V|$ e $\alpha$ é o coeficiente de \textit{clustering} local desse vértice. Nesse caso $\epsilon$ e $\delta$ são tratados como constantes \cite[677--686]{Kutzkov2013}. Temos também a abordagem de contagem
de triângulos locais em algoritmos de aproximação como o apresentado por Beccheti et al. em
2010: um algoritmo de \textit{semi-streaming} para calcular uma aproximação com erro relativo do número de triângulos de cada vértice de um grafo em tempo $O(m \cdot k)$, onde $k$ é o número de
passagens sobres os dados no algoritmo \cite[1--28]{Becchetti2010}. Uma excelente revisão sobre esses métodos, incluindo algoritmos exatos e de aproximação, para o problema de contagem global de triângulos foi feita por Al Hasan e Dave em 2018 \cite{AlHasan2018}.

Uma abordagem que pode ser utilizada para algoritmos de aproximação é a teoria de
complexidade de amostra para se obter melhores resultados. O algoritmo proposto por Lima et al. (2022) – que é o objeto de estudo deste trabalho, tendo como objetivo implementá-lo – faz uso da \textit{VC-Dimension} e \textit{$\epsilon$-sample}, onde amostra arestas de um grafo de entrada $G$ e, para
constantes fixas $0 < \epsilon,\delta,p < 1$, fornece uma estimativa $l'(v)$ para o valor exato $l(v)$ do coeficiente de \textit{clustering} local de cada vértice $v \in V$. Os resultados respeitam $|l(v)-l'(v)| \leq \epsilon l(v)$, com probabilidade de pelo menos $1 - \delta$ sempre que $l(v)$ for pelo menos $\frac{pm}{2^{\delta_v}}$, onde $\delta_v$ é o grau de $v$ \cite[64--65]{Lima2022}. A utilização da \textit{VC-Dimension} e \textit{$\epsilon$-sample} ao invés de abordagens clássicas da literatura, como o limitante da união, para algoritmos de amostragem é em decorrência da importância que a estrutura combinatória das entradas se faz presente neles.

Existe também um algoritmo proposto por Zhang et al. em 2015 que busca computar os top-$k$ vértices com maiores coeficientes de \textit{clustering} local de um grafo e utiliza \textit{VC-Dimension} e o teorema \textit{$\epsilon$-sample} \cite[1445--1454]{Zhang2015}, porém em um contexto e espaço de intervalos diferentes do proposto por Lima et al. (2022).